\documentclass[12pt]{article}
\usepackage{amsmath, amssymb}

\title{CSE 400: Fundamentals of Probability in Computing \\ 
Lecture 25 -- Tutorial Session: Problem Solving (Inequalities)}
\date{April 9, 2026}

\begin{document}

\maketitle

\section*{Q1: Markov’s Inequality}

\textbf{Given:} \\
Average annual income = \$40,000

\textbf{(a) Requirement:} \\
Use Markov’s inequality to find an upper bound on:
\[
P(\text{income} \geq 120{,}000)
\]

\textbf{(b) Requirement:} \\
Determine whether Markov’s inequality can be used to bound:
\[
P(\text{income} < 10{,}000)
\]
Justify the answer.

\section*{Q2: Chebyshev’s Inequality}

\textbf{Given:} \\
Mean: $\mu = 100\,\Omega$ \\
Standard deviation: $\sigma = 2\,\Omega$ \\
Defective range: outside $[94\,\Omega, 106\,\Omega]$

\textbf{(a) Requirement:} \\
Use Chebyshev’s inequality to find an upper bound on the probability that a resistor is defective.

\textbf{(b) Requirement:} \\
If 10,000 resistors are produced, determine the maximum expected number of defective units.

\section*{Q3: Convexity and Expectation Inequality}

\textbf{Given:} \\
Random variable $Y$ such that $0 \leq Y \leq 1$ \\
$E[Y] = \mu$ \\
Function: $g(x) = e^x$

\textbf{Requirement:} \\
Using convexity of $g(x)$, prove:
\[
E[e^{tY}] \leq 1 - \mu + \mu e^t
\]

\section*{Q4: Jensen’s Inequality}

\textbf{Given:} \\
Random variable $X > 0$ \\
$E[X] = \mu$ \\
Cost function: $C(X) = X^a$, where $a > 1$

\textbf{Requirement:} \\
Using Jensen’s inequality, prove:
\[
E[X^a] \geq (E[X])^a
\]

\section*{Q5: Chebyshev’s Inequality (Standard Form)}

\textbf{Given:} \\
Random variable $X$ with mean $\mu$ \\
Variance $\sigma^2$ \\
Real number $k > 0$

\textbf{Requirement:} \\
Show that:
\[
P(|X - \mu| \geq k\sigma) \leq \frac{1}{k^2}
\]

\section*{Q6: Production Random Variable}

\textbf{Given:} \\
Weekly production is a random variable \\
Mean = 500

\textbf{(a) Requirement:} \\
Determine what can be said about:
\[
P(\text{production} \geq 1000)
\]

\textbf{(b) Additional Given:} \\
Variance = 100

\textbf{Requirement:} \\
Determine what can be said about:
\[
P(400 \leq \text{production} \leq 600)
\]

\section*{Q7: Chernoff Bounds (Independent Random Variables)}

\textbf{Given:} \\
Independent trials $X_i \in \{+1, -1\}$ \\
$P(X_i = 1) = P(X_i = -1) = \frac{1}{2}$

\[
S_n = \sum_{i=1}^{n} X_i
\]

\textbf{(a) Requirement:} \\
Show:
\[
E[e^{tX_i}] \leq e^{t^2/2}
\]

\textbf{(b) Requirement:} \\
Deduce:
\[
E[e^{tS_n}] \leq e^{nt^2/2}
\]

\section*{Q8: Pairing Problem}

\textbf{Given:} \\
200 people: 100 men and 100 women \\
Randomly divided into 100 pairs

\textbf{Requirement:} \\
Find an upper bound on the probability that at most 30 pairs consist of one man and one woman.

\section*{Q9: Inequality Proof}

\textbf{Requirement:} \\
Show that:
\[
1 + \frac{1}{2^2} + \frac{1}{3^2} + \cdots + \frac{1}{n^2} \geq \frac{2^n}{n + 2^n + 5}
\]

\section*{Q10: Poisson Random Variable}

\textbf{Given:} \\
$X \sim \text{Poisson}(20)$

\textbf{(a) Requirement:} \\
Use Markov’s inequality to bound:
\[
p = P(X \geq 26)
\]

\textbf{(b) Requirement:} \\
Use one-sided Chebyshev inequality to bound $p$.

\textbf{(c) Requirement:} \\
Use Chernoff bound to bound $p$.

\end{document}